\chapter{Introduction and Overview}

\begin{flushright}
\textit{Sensitivity analysis begins when a model is forced to say what matters.}
\end{flushright}

\bigskip

An infectious disease has just entered a new country.
The public health response team has built a compartmental model with fourteen
parameters: transmission rates, recovery rates, contact patterns,
immunity duration, disease-induced mortality, and more.
The team has funding to deploy exactly two interventions
and time to measure exactly three parameters with high precision
before a policy decision must be made.

The question due by tomorrow morning is this: \emph{which three parameters
are worth measuring, and which two interventions will most effectively
reduce the spread?}

This is the problem that sensitivity analysis is designed to answer.
Before the mathematics, we need a language for asking the question precisely.

%%------------------------------------------------------------
\section{The Central Question}
\label{sec:centralquestion}
%%------------------------------------------------------------

Every mathematical model takes inputs and produces outputs.
The inputs include parameters whose exact values may be uncertain
or whose values might be changed through intervention.
The outputs include whatever the model is designed to predict.
Sensitivity analysis asks: \emph{which inputs have the most
effect on the outputs we care about?}

To answer this precisely, we adopt two terms used throughout the book.

\medskip

A \textbf{Parameter of Interest} (\POI)\index{parameter of interest (POI)} is any input to the model
whose variation is worth studying. In practice, a parameter becomes a \POI
for one of three reasons, and the reason matters for how you interpret
the sensitivity index.

\begin{enumerate}

\item \textbf{Epistemic uncertainty} (uncertainty from lack of knowledge).
The parameter is fixed by nature but imperfectly known: different studies
report different estimates, the data are sparse, or the measurement process
is imprecise. Epistemic uncertainty is \emph{reducible} --- better data,
improved experiments, or longer observation windows can sharpen the estimate.
Sensitivity analysis identifies which epistemically uncertain parameters
most affect the output, directing limited data-collection resources
toward the measurements that would reduce output uncertainty the most.
In the disease example, the transmission probability per contact is
epistemically uncertain: studies estimate it from contact-tracing data,
but estimates vary across populations and settings.

\item \textbf{Aleatory uncertainty} (inherent stochastic variability).
The parameter genuinely varies from individual to individual, place to
place, or time to time, and this variation is \emph{irreducible}---
it persists even with unlimited data. More observations characterize
the distribution more precisely but do not eliminate the randomness.
In the disease example, the contact rate varies stochastically
across the population: individuals differ in how many contacts they
make per day. Sensitivity analysis with aleatory uncertainty typically
asks how much of the output variance is attributable to each
stochastic source, motivating the global SA methods of Chapter~9.

\item \textbf{Control variable} (a lever the investigator can move).
The parameter is known and could, in principle, be changed by a policy
decision, intervention, or design choice. Sensitivity analysis in this
context answers the question: which intervention has the most leverage
over the outcome I care about? In the disease example, the contact rate
can be reduced by social distancing policies, and the transmission
probability can be reduced by mask requirements or vaccination. The
sensitivity index tells the policy maker which reduction --- contacts
or transmission --- produces the largest reduction in $R_0$.

\end{enumerate}

Many parameters fall into more than one category simultaneously:
the contact rate is both uncertain (its current value is imprecisely
known) and controllable (it can be targeted by policy).
In practice, the interpretation of a sensitivity index should
always be read in light of which role the \POI is playing.
A large sensitivity index for an epistemically uncertain parameter
is a call to collect better data; for an aleatory parameter, it is
a call to characterize the distribution more carefully; for a control
variable, it is a call to prioritize that intervention.

In the disease example above, the transmission probability
per contact is primarily a \POI because it is epistemically uncertain
(different studies report different values) and because it can be
targeted by an intervention (vaccines reduce it).
The contact rate is a \POI for both the same reasons:
it varies across populations (epistemic and aleatory uncertainty)
and can be reduced through social distancing policies (control variable).

A \textbf{Quantity of Interest} (\QOI)\index{quantity of interest (QOI)} is any model output that the
investigator cares about predicting or controlling. The peak number of
infectious individuals, the total number of deaths over a six-month
period, and the basic reproduction number $\mathcal{R}_0$ are all
plausible \QOIs for an epidemic model. The choice of \QOI is not neutral:
different choices lead to different rankings of parameter importance.

The relationship between \POIs and \QOIs is the subject of this book.
Specifically, we want to know how much each \POI affects each \QOI,
so that we can rank parameters by importance, design interventions
efficiently, and allocate measurement resources wisely.

One of the most persistent misconceptions in mathematical modeling
is the assumption that all parameters matter equally until proven
otherwise. They do not. In virtually every model studied in this
book, a small number of \POIs account for most of the variation
in the \QOIs, while the remaining parameters have negligible effect.
Identifying that small number is one of the primary goals of
sensitivity analysis, and the answer is rarely obvious from the
model equations alone.

\begin{selfcheck}
In the outbreak scenario above, classify each of the following
as a \POI, a \QOI, or neither:
the transmission probability per contact;
the peak number of infectious individuals;
the day the first case was reported;
the mean number of days a patient is infectious.
\end{selfcheck}
\begin{scAnswer}
\POI: transmission probability (uncertain; can be reduced by vaccines),
mean infectious period (uncertain; can be reduced by treatment).
\QOI: peak number of infectious individuals (model output we care about).
Neither: day of first case (a historical fact, not a model parameter
or prediction). Note that the mean infectious period could also be
a \QOI in a different study; this classification depends on
the question being asked.
\end{scAnswer}

Figure~\ref{fig:forwardproblem} illustrates the basic structure.
The forward problem takes nominal \POI values as input
and produces the \QOI as output. Sensitivity analysis adds one
more step: it asks how the \QOI changes as the \POIs vary.

\begin{figure}[htb]
  \begin{center}
  \begin{tikzpicture}[
    node distance = 2.6cm,
    box/.style = {
      rectangle, rounded corners=4pt, draw=black, thick,
      minimum width=2.5cm, minimum height=1.1cm,
      text centered, font=\small
    },
    sarrow/.style = {-{Latex[length=3mm]}, thick}
  ]
    %% Main flow nodes — single-line labels avoid \\[...] issues
    \node[box, fill=black!8]  (pois)  {\textbf{\POIs}};
    \node[box, fill=black!15, right=of pois]  (model) {\textbf{Model}};
    \node[box, fill=black!8,  right=of model] (qois)  {\textbf{\QOIs}};

    %% Sub-labels below each box
    \node[font=\footnotesize\itshape, below=2pt of pois]  {parameters};
    \node[font=\footnotesize\itshape, below=2pt of model] {equations};
    \node[font=\footnotesize\itshape, below=2pt of qois]  {outputs};

    %% Forward arrows
    \draw[sarrow] (pois)  -- node[above, font=\footnotesize]{$\vec{p}$} (model);
    \draw[sarrow] (model) -- node[above, font=\footnotesize]{$\vec{q}(\vec{p})$} (qois);

    %% SA brace and formula below
    \draw[thick, decorate,
          decoration={brace, amplitude=6pt, mirror}]
      ([yshift=-18pt]pois.south west) -- ([yshift=-18pt]qois.south east)
      node[midway, below=9pt, font=\small\itshape]{
        Sensitivity Analysis:\quad
        $\SI{p_i}{q_j} = \dfrac{p_i}{q_j}\,\dfrac{\partial q_j}{\partial p_i}$
      };
  \end{tikzpicture}
  \end{center}
  \caption{The forward problem takes nominal values for the input \POIs
           and produces the associated output \QOI. Sensitivity analysis
           quantifies how changes in the \POIs propagate to the \QOI
           by computing the normalized partial derivative $\SI{p_i}{q_j}$
           for each \POI--\QOI pair.}
  \label{fig:forwardproblem}
\end{figure}

%%------------------------------------------------------------
\section{Two Directions: Forward and Adjoint Sensitivity}
\label{sec:twoDirections}
%%------------------------------------------------------------

There are two complementary ways to ask the sensitivity question,
and they differ in a fundamental way.

Figure~\ref{fig:forwardadjoint} captures the distinction geometrically.
In panel~(a), a perturbation originates at a specific input location
and spreads forward through the model, influencing a broad region of
the output. This is \textbf{forward sensitivity analysis}\index{forward sensitivity analysis}.
Panel~(b) reverses the direction: a specified variation in the output
is traced backward to the input locations that could have caused it.
This is \textbf{adjoint sensitivity analysis}\index{adjoint sensitivity analysis}.

\begin{figure}[htb]
  \begin{center}
  \begin{tikzpicture}[
    modelbox/.style = {
      rectangle, rounded corners=5pt, draw=black, line width=1.2pt,
      minimum width=1.6cm, minimum height=4.4cm,
      text centered, fill=black!12, font=\small\bfseries
    },
    qoibox/.style = {
      rectangle, rounded corners=4pt, draw=black, thick,
      minimum width=1.4cm, minimum height=0.9cm,
      text centered, fill=black!6, font=\small
    },
    farrow/.style  = {-{Latex[length=3mm]}, thick},
    garrow/.style  = {-{Latex[length=2.5mm]}, thin, black!35},
    aarrow/.style  = {-{Latex[length=3.5mm]}, very thick, dashed},
    seclabel/.style = {font=\small\bfseries},
    sublabel/.style = {font=\footnotesize\itshape},
    plab/.style     = {font=\small},
  ]

    %%----------------------------------------------------------
    %% Panel (a): Forward SA
    %% One POI (p_j) is perturbed; its effect flows through the
    %% Model and appears at the QOI. Other POIs shown faded to
    %% indicate the process repeats for each one.
    %%----------------------------------------------------------
    \node[seclabel] at (2.8, 2.55) {(a) Forward SA \quad ``What if?''};

    %% POI column
    \node[plab]                      (fp1) at (0.0,  1.50) {$p_1$};
    \node[plab]                      (fp2) at (0.0,  0.75) {$p_2$};
    \node[font=\small]                     at (0.0,  0.10) {$\vdots$};
    \node[plab, font=\small\bfseries](fpj) at (0.0, -0.55) {$p_j$};
    \node[font=\small]                     at (0.0, -1.10) {$\vdots$};
    \node[plab]                      (fpm) at (0.0, -1.70) {$p_m$};

    \node[modelbox] (fmod) at (2.8, -0.10) {Model};
    \node[qoibox]   (fqoi) at (5.4, -0.10) {\QOI};

    %% Faded arrows: all other POIs enter model at their y-level
    \draw[garrow] (fp1.east) -- (fmod.west |- fp1.center);
    \draw[garrow] (fp2.east) -- (fmod.west |- fp2.center);
    \draw[garrow] (fpm.east) -- (fmod.west |- fpm.center);

    %% Bold path: p_j through Model to QOI
    \draw[farrow, line width=1.4pt] (fpj.east) -- (fmod.west |- fpj.center);
    \draw[farrow, line width=1.4pt] (fmod.east) -- (fqoi.west);

    %% Perturbation indicator at p_j
    \fill[black!70] (-0.85,-0.55) circle (4pt);
    \draw[->,thick] (-0.85,-0.55) -- (fpj.west);
    \node[sublabel, left=20pt of fpj, align=right] {perturb\\$p_j$};

    \node[sublabel, align=center, text width=5.0cm] at (2.8, -3.20)
      {Perturb one \POI, observe the \QOI.\\
       Repeat for each \POI: $m$ solves total.};

    %%----------------------------------------------------------
    %% Dividing rule
    %%----------------------------------------------------------
    \draw[thin, gray!50] (7.0,-3.50) -- (7.0,2.80);

    %%----------------------------------------------------------
    %% Panel (b): Adjoint SA
    %% A specified QOI variation enters the Model from the right
    %% and flows BACKWARD, producing all POI sensitivities at once.
    %% Arrowhead at END of each dashed path (−{Latex}): points left.
    %%----------------------------------------------------------
    \node[seclabel] at (10.5, 2.55) {(b) Adjoint SA \quad ``How come?''};

    %% POI column
    \node[plab] (ap1) at (7.7,  1.50) {$p_1$};
    \node[plab] (ap2) at (7.7,  0.75) {$p_2$};
    \node[font=\small]  at (7.7,  0.10) {$\vdots$};
    \node[plab] (apj) at (7.7, -0.55) {$p_j$};
    \node[font=\small]  at (7.7, -1.10) {$\vdots$};
    \node[plab] (apm) at (7.7, -1.70) {$p_m$};

    \node[modelbox] (amod) at (10.5, -0.10) {Model};
    \node[qoibox]   (aqoi) at (13.1, -0.10) {\QOI};

    %% ONE dashed arrow: QOI → Model (backward; arrowhead enters model)
    \draw[aarrow] (aqoi.west) -- (amod.east);

    %% ALL dashed arrows: Model → each POI (backward; arrowhead at each POI)
    \draw[aarrow] (amod.west |- ap1.center) -- (ap1.east);
    \draw[aarrow] (amod.west |- ap2.center) -- (ap2.east);
    \draw[aarrow] (amod.west |- apj.center) -- (apj.east);
    \draw[aarrow] (amod.west |- apm.center) -- (apm.east);

    %% "specify δq_j" below the QOI box
    \node[sublabel, below=5pt of aqoi, align=center] {specify $\delta q_j$};

    \node[sublabel, align=center, text width=5.2cm] at (10.5, -3.20)
      {Specify one \QOI variation, solve adjoint backward.\\
       All \POI sensitivities recovered: 1 solve.};

  \end{tikzpicture}
  \end{center}
  \caption{(a)~Forward sensitivity analysis answers ``What if?''\ questions.
           A perturbation to one input \POI $p_j$ propagates forward through
           the model and produces a change in the output \QOI\@.
           Because the model must be solved once for each \POI, forward
           sensitivity requires $m$ solves for $m$ \POIs.
           (b)~Adjoint sensitivity analysis answers ``How come?''\ questions.
           A specified variation $\delta q_j$ in the output \QOI enters the
           model from the right and is propagated \emph{backward} through it.
           A single adjoint solve simultaneously recovers the sensitivity of
           $q_j$ with respect to every \POI $p_1, \ldots, p_m$.
           The adjoint is developed in Chapters~5 and~6; for now, the key
           contrast is computational: forward SA costs $m$ solves,
           adjoint SA costs~$1$.}
  \label{fig:forwardadjoint}
\end{figure}

The ``What if?'' and ``How come?'' labels are not merely rhetorical.
They correspond to genuinely different computational strategies
and genuinely different questions.

\subsection{Forward Sensitivity Analysis: ``What if?''}

Consider a disease model for an emerging outbreak.
The transmission probability $\beta$ is estimated from early case data,
but the estimate carries significant uncertainty. A public health officer
asks: \emph{if our estimate of $\beta$ is wrong by 10\%, how much
does our predicted peak infection count change?}

This is a forward question. We are tracing the effect of a change
in an input \POI (the transmission probability) forward to an
output \QOI (the peak infection count). Forward sensitivity analysis
answers this question by computing a sensitivity index: a number that
tells us, for a small percentage change in a \POI, what percentage
change to expect in the \QOI.

A sensitivity index of $+0.8$ would mean that a 10\%\ increase in
$\beta$ produces approximately an $8\%$ increase in the peak infection
count. A sensitivity index of $-1.0$ for the recovery rate would mean
that a 10\%\ increase in the recovery rate produces a 10\%\ decrease
in the peak. The sign tells you the direction; the magnitude tells you
the amplification. Chapter~2 develops this definition carefully and
shows how to compute these indices for a wide range of models.

\begin{selfcheck}
A forward sensitivity analysis gives $S_{\tau}^{R_0} = +1.2$ (the notation is defined formally in Chapter~2),
where $\tau$ is the mean infectious period (in days) and $R_0$ is
the basic reproduction number.
If the mean infectious period increases by $5\%$, approximately
what happens to $R_0$? Is this good news or bad news for disease control?
\end{selfcheck}
\begin{scAnswer}
$R_0$ increases by approximately $5\% \times 1.2 = 6\%$. This is bad news:
a larger $R_0$ means more transmission. The positive sign tells us that
\POI and \QOI move in the same direction, so longer infectious periods
lead to more disease spread.
\end{scAnswer}

\subsection{Adjoint Sensitivity Analysis: ``How come?''}

Now consider a different problem. City planners want to site a new
industrial facility near a residential area. The facility emits
pollutants, and regulators have set a standard: the average pollutant
concentration at a specified monitoring station must not exceed a
given threshold. A transport model predicts whether any candidate
site meets this standard.

The planner's question is different in character from the previous one.
Rather than asking ``if this parameter changes, what happens to the
output?'', they want to know: \emph{for a given allowed change in the
output (the pollutant concentration), how much variation in the inputs
(wind speed, emission rate, atmospheric stability) is permissible?}

This is an adjoint question. It runs the sensitivity analysis in
reverse: from a specified variation in the \QOI back to the
\POIs that could have caused it. For models with many \POIs
and only a few \QOIs, adjoint sensitivity analysis is
dramatically more efficient than computing forward sensitivities
for every parameter individually. Chapters~5 and 6 explain
why this is so and how to exploit it.

\begin{selfcheck}
An environmental regulator measures pollutant concentration at a
monitoring station and finds it exceeds the legal threshold.
She wants to determine which weather assumptions in the
transport model are most likely responsible for the exceedance.
Is this a forward or adjoint question? Explain in one sentence.
\end{selfcheck}
\begin{scAnswer}
Adjoint. The regulator starts with a specified change in the \QOI
(exceedance of the threshold) and traces backward to identify which
input \POIs (weather assumptions) most plausibly caused it.
A forward analysis would instead ask how changing a specific weather
assumption would affect the predicted concentration, which runs in the opposite
direction.
\end{scAnswer}

%%------------------------------------------------------------
\section{Global Sensitivity: When Local Is Not Enough}
\label{sec:globalSA}
%%------------------------------------------------------------

Both forward and adjoint sensitivity analysis are \emph{local}:
they measure the effect of small perturbations to the \POIs
near a fixed nominal value. This is powerful, but it assumes
that the model behaves approximately linearly in the vicinity
of that nominal value.

Many real models are not linear, and many \POIs carry substantial
uncertainty that cannot be called ``small.'' In these cases,
local sensitivity indices can be misleading. A \POI with a
small local sensitivity index may still have a large effect
on the \QOI when varied across its full range of plausible values,
if the model is nonlinear in that parameter.

\begin{example}[Local SI does not capture global behavior]
\label{ex:local_global}
Consider $q(p) = p^2$ and suppose the nominal value is $p^* = 0.1$.
The local sensitivity index is $S_p^q = (p/q)\,(dq/dp) = (p/p^2)(2p) = 2$,
which looks moderate.
But if $p$ is uncertain and could range from $0.1$ to $1.0$,
the QOI ranges from $0.01$ to $1.0$ --- a hundredfold change.
A local SI computed at $p^* = 0.1$ gives no information about
this hundredfold range.
Global SA, introduced in Chapter~9, is designed to capture exactly
this kind of large-range nonlinear behavior.
\end{example}

Chapter~9 introduces \textbf{global sensitivity analysis}\index{global sensitivity analysis},
which asks how the output varies when the \POIs are varied
across their entire plausible ranges simultaneously.
The primary tool there, the Sobol index, measures
how much of the total output variance is attributable to
each \POI. Global SA does not replace local SA; it
complements it, answering questions that local SA cannot.

%%------------------------------------------------------------
\section{Sensitivity Analysis and Uncertainty Quantification}
\label{sec:sauq}
%%------------------------------------------------------------

Sensitivity analysis and uncertainty quantification are
related but distinct enterprises.

Sensitivity analysis, as developed in this book, is
\emph{deterministic}: we assign nominal values to the \POIs
and compute how the \QOI changes when those values are
perturbed. The \POIs are treated as numbers with specific
values, not as random variables. The sensitivity index
is a derivative: a precise mathematical object that exists
when the model is differentiable.

Uncertainty quantification treats the \POIs as random variables
with associated probability distributions. The goal is to
determine the distribution of the \QOI given the distributions
of the \POIs. This is a statistical enterprise, and its primary
tools (Monte Carlo methods, Bayesian inference, polynomial chaos)
are the subject of the companion volume by Smith~\cite{smith2014uncertainty}.
Smith's Chapters~9 and~10 develop local and global sensitivity analysis
at the graduate level, extending the material in this book with
functional-analytic and measure-theoretic foundations.
Smith's Chapter~6 treats parameter selection using the sensitivity
Jacobian introduced in our Chapter~4, connecting it to the
surrogate-model literature.
The SA foundations developed here follow Arriola and Hyman~\cite{arriola2007being,arriola2009sensitivity}.
For the broader development of SA as a formal discipline, see Saltelli et al.~\cite{saltelli2021sensitivity},
Razavi et al.~\cite{razavi2021future}, and the annotated historical timeline in
Tarantola et al.~\cite{tarantola2024annotated}.

The connection between the two fields is real and important.
Global SA, covered in Chapter~9, bridges the gap: Sobol indices
are defined in terms of variances, which requires probabilistic
inputs, but they can be estimated by methods that build directly
on the ideas in this book. Students who complete this course
will find that the second semester's UQ material makes more
sense because of the SA foundation established here.

%%------------------------------------------------------------
\section{Three Ideas That Run Through the Book}
\label{sec:threeIdeas}
%%------------------------------------------------------------

Every chapter of this book revisits one or more of three
fundamental ideas. We name them here so that students can
recognize them when they reappear, each time in a new setting.

\medskip

\noindent\textbf{Unifying Idea~1:\index{unifying ideas} The sensitivity index is a
normalized derivative.}

At its core, SA answers a single question: if this input changes
by a small percentage, by what percentage does the output change?
That ratio is the sensitivity index. It is not a mysterious new
object: it is the derivative, normalized so that its magnitude
is independent of units and scales.

This idea appears first in Chapter~2, where we define the
sensitivity index formally. It reappears in Chapter~3 as a
finite-difference approximation, in Chapters~4--6 as
an analytic derivative, and in Chapter~9 as the foundation
for Sobol indices. Whenever a sensitivity quantity appears
in this book, it is an instance of this one idea.
Smith~\cite{smith2014uncertainty} develops the same normalized-derivative
framework in \S9.1--9.2 and connects it to experimental design.

\medskip

\noindent\textbf{Unifying Idea~2: The adjoint is the transpose
in whatever space you are working in.}

The adjoint sensitivity method, which makes it possible
to compute the sensitivities of one \QOI with respect to
hundreds of \POIs at the cost of solving just two problems
instead of hundreds, rests on a single mathematical operation:
transposition. For a linear system $A\vec{u} = \vec{b}$,
the adjoint involves the transposed matrix $A^T$.
For a computation graph (the structure underlying every
computer program), the adjoint is the chain rule applied in
reverse order, which is the backpropagation operation familiar from
machine learning. For an ordinary differential equation,
the adjoint is a second ODE running backward in time.

These three forms look different on the surface.
They are the same idea.

Chapter~5 introduces the adjoint through the matrix transpose,
which is the simplest and most familiar form. Chapter~6 extends
it to differential equations. At each transition, the book
explicitly identifies the connection back to this unifying idea.

\medskip

\noindent\textbf{Unifying Idea~3: Local SA is a linearization.}

Every sensitivity index in Chapters~2--8 is based on a derivative,
and a derivative is a linear approximation to a generally nonlinear
function. This approximation is valid when the perturbation is small
and the model is smooth. It fails near bifurcation points, where
the model changes qualitative behavior. It fails when the \POI
uncertainty is large. It fails when the model is strongly nonlinear
in the \POI of interest.

Chapter~8 (\emph{Caveat Emptor}) documents the conditions
under which local SA breaks down. Chapter~9 provides
the global alternative. Both chapters are applications
of this one principle: local SA is a linearization,
and linearizations have limits.

%%------------------------------------------------------------
\section{A Note on Parameterization}
\label{sec:parameterization}
%%------------------------------------------------------------

One practical decision that profoundly affects the usefulness
of SA results deserves mention before Chapter~2.

The sensitivity index depends on how a \POI is defined.
In a compartmental epidemic model, the recovery process can
be parameterized either by the \emph{recovery rate} $\gamma$
(units: per day) or by the \emph{mean infectious period}
$\tau = 1/\gamma$ (units: days). These two parameterizations
are mathematically equivalent, each determining the other.
But their sensitivity indices are not equal: they have
the same magnitude and opposite signs.

Which one should you use? The answer is: whichever one you can
discuss in ordinary language without confusion. People who work
with epidemic models (clinicians, public health officials,
and policy makers) think naturally in terms of days.
A clinician understands immediately what it means to say
``the mean infectious period is 7 days.'' They do not
think in recovery rates. Reporting $\SI{\gamma}{R_0} = -1.0$
requires the reader to invert $\gamma$ mentally before the
result makes intuitive sense. Reporting $\SI{\tau}{R_0} = +1.0$
is immediately clear: a longer infectious period increases
transmission.

Throughout this book, we choose \POIs that can be understood
in everyday terms: durations rather than rates, probabilities
rather than their complements, counts rather than densities
when the distinction matters. This is not just a matter of
aesthetics. Sensitivity indices that require mental inversion
to interpret are sensitivity indices that will be misread.

Chapter~2 formalizes this principle and shows how to convert
between parameterizations. Readers who find that their model's
natural parameters are rates will learn how to restate
their results in terms of the corresponding durations
before reporting them.

%%------------------------------------------------------------
\section{Structure of the Book}
\label{sec:structure}
%%------------------------------------------------------------

The book is organized in five parts. Each part builds on
the previous, and the SIR epidemic model~\cite{kermack1927contribution,wu2013sensitivity} serves as the
recurring example throughout, so that students can track
a single application as the mathematical tools become
progressively more powerful.

\medskip

\textbf{Part~I: The Sensitivity Question} (Chapters~1--2)
establishes the central problem and introduces the sensitivity
index as the primary tool. Chapter~2 develops the formal
definition, the sign interpretation, the tornado plot
as a display format, and the principle of interpretable
parameterization. By the end of Part~I, students can compute
and interpret sensitivity indices for any differentiable
function of several variables.

\textbf{Part~II: Computing Sensitivity in Practice}
(Chapter~3) gives students an immediately usable computational
method for any model, including models that cannot be
differentiated analytically. Finite-difference approximations,
step-size selection, the sensitivity Jacobian, and the
treatment of correlated \POIs are developed here.
Students who work through Chapter~3 can perform SA on any
model they can run, including inherited code and black-box
simulation software.

\textbf{Part~III: Analytic Forward Sensitivity Analysis}
(Chapter~4) shows how to differentiate model equations directly
to obtain exact sensitivity information at lower computational
cost. The forward sensitivity equation is derived for algebraic
systems, difference equations, and ordinary differential
equations. The chapter ends by quantifying the computational
bottleneck that motivates the adjoint approach.

\textbf{Part~IV: The Adjoint Approach} (Chapters~5--7)
introduces adjoint sensitivity analysis through three
progressively more abstract settings: linear systems
in Chapter~5, dynamical systems in Chapter~6, and
computational tools including algorithmic differentiation
in Chapter~7. The central theme is Unifying Idea~2:
the adjoint is the matrix transpose in whatever
mathematical space the problem inhabits.

\textbf{Part~V: Limitations and Extensions}
(Chapters~8--9) addresses the boundaries of local SA.
Chapter~8 documents the conditions under which sensitivity
indices give misleading results: bifurcation points,
ill-conditioned systems, and method-question mismatch.
Chapter~9 introduces global SA through variance decomposition,
Sobol indices, and partial rank correlation coefficients.
It closes with an explicit bridge to the second semester,
where Smith's \emph{Uncertainty Quantification}~\cite{smith2014uncertainty}
takes over from the point where this book ends.

\medskip

\noindent\textbf{Computing environment.} Every chapter
includes a structured MATLAB laboratory that implements
the chapter's core idea on the SIR epidemic model.
Students who complete these laboratories will leave the
course with a portable toolkit for performing SA on any
model they encounter in research or practice.
The complete MATLAB code library --- production-quality functions,
automated test suites, and usage examples --- is available at:
\begin{center}
  \texttt{https://github.com/mhyman/arriola2026sa}
\end{center}
Each function in the repository corresponds directly to a concept
or algorithm in the text. \texttt{sir\_nominal.m} defines the SIR
parameters used in every chapter; \texttt{sensitivity\_jacobian.m}
implements the centered-difference Jacobian from Chapter~3;
\texttt{sobol\_jansen.m} implements the Jansen estimators from
Chapter~9. The \texttt{src/} directory is organized to mirror
the book's structure, so students can find the code for any
chapter without searching.

\noindent\textbf{Notation.} Throughout the book, the
sensitivity index of a \QOI~$q$ with respect to a
\POI~$p$ is written $\SI{p}{q}$.
\POIs are always chosen to be interpretable in everyday
terms: durations, probabilities, counts. The relationship
to the mathematical parameters in the model equations
is made explicit wherever the two differ.

\noindent\textbf{Scope note: scalar QOIs.}
Every sensitivity index in this book is computed for a
\emph{scalar} QOI --- a single number such as the peak infection
count, the total disease burden, or the value of $R_0$.
Real applications sometimes require sensitivity of an entire output
\emph{trajectory}: how does the full time series $I(t)$ change when
$\beta$ changes?
Extending SA to functional (trajectory-valued) QOIs requires either
reducing the trajectory to a scalar (which loses information) or
computing a time-dependent SI at each time point (as in Chapter~4's
Figure~\ref{fig:sir_timevarySI}).
A rigorous treatment of functional sensitivity, using the adjoint ODE
to compute sensitivity with respect to the entire output path, is
developed in Smith~\cite{smith2014uncertainty} \S9.4 and is a natural
next step after mastering the scalar case.

%%------------------------------------------------------------
\section{Exercises}
%%------------------------------------------------------------

Each exercise carries a Bloom's Taxonomy tag indicating cognitive level:
\textit{L2} (understand), \textit{L3} (apply), \textit{L4} (analyze),
\textit{L5} (evaluate), \textit{L6} (create and synthesize).
Higher-level exercises require constructing an argument or designing a procedure,
not just applying a formula.

\begin{exercise}[\bloom{L2}]
\label{ex:ch1:vocabulary}
Explain in your own words the difference between a \POI and a \QOI.
For a model describing the spread of antibiotic-resistant bacteria
in a hospital, give two examples of each. Then explain, in one sentence,
what sensitivity analysis would tell you about this model that
simply running the model would not.
\end{exercise}

\begin{exercise}[\bloom{L2}]
\label{ex:ch1:forwardadjoint}
What is the difference between forward and adjoint sensitivity analysis?
Describe each in terms of the ``What if?'' and ``How come?'' framing.
For each of the following questions, classify it as a \emph{forward question}
or an \emph{adjoint question} based on the type of question being asked
(not on which algorithm would be most efficient to use computationally ---
efficiency is the subject of Chapters~5 and~6).
\begin{enumerate}[label=(\alph*)]
  \item If the mean latent period increases by $10\%$, how does the peak
        infection count change?
  \item The peak infection count is higher than expected. Which parameter
        uncertainties are most likely to have caused this?
  \item Among all parameters in the model, which one has the largest
        effect on the basic reproduction number?
\end{enumerate}
\end{exercise}

\begin{exercise}[\bloom{L4}]
\label{ex:ch1:methodchoice}
For each of the following modeling scenarios, state whether forward SA,
adjoint SA, or global SA is most appropriate. Justify each choice in
one or two sentences.
\begin{enumerate}[label=(\alph*)]
  \item A climate model has 500 tunable parameters.
        The investigator wants to know which 10 parameters most
        affect the mean surface temperature over a 100-year run.
        Each model run takes 8 hours.
  \item A pollution transport model has 6 input parameters.
        The regulator wants to know: if the pollution at a specific
        monitoring station exceeds the legal threshold,
        which input assumptions most likely caused this?
  \item A pharmacokinetic model has 4 parameters, each of which
        is known only to within a factor of 2. The investigator
        wants to know the range of predicted drug concentrations
        consistent with this uncertainty.
\end{enumerate}
\end{exercise}

\begin{exercise}[\bloom{L5}]
\label{ex:ch1:studydesign}
You are advising a team modeling the transmission of dengue fever.
The model has 9 parameters.
The team has funding to measure exactly 3 parameters precisely in
the field and to implement exactly 2 vector-control interventions.
Write a one-page plan describing how you would use SA to guide
these decisions. Your plan should specify:
what \QOI you would focus on and why;
how you would determine which 3 parameters to measure;
what additional information (if any) you would need
beyond the SA results to recommend interventions;
and one important limitation of your approach
that the team should be aware of.
\end{exercise}

\begin{exercise}[\bloom{L6}]
\label{ex:ch1:evaluation}
A published report on a disease model states:
``Sensitivity analysis was performed on all 22 model parameters
simultaneously, demonstrating that the model is robust
to parameter uncertainty.'' The analysis was conducted by
computing local sensitivity indices at the best-fit parameter values.

Identify at least two substantive problems with this description
or its conclusions. For each problem you identify, explain
specifically what the authors should have done differently.
\end{exercise}

%%------------------------------------------------------------
%% Bibliography for standalone compilation
%%------------------------------------------------------------

